% ═══════════════════════════════════════════════════════════════════════════
% LEHRERHANDREICHUNG: DS-10 Variante B (Speed-Dating)
% Kompakte Übersicht auf 1 Seite
% ═══════════════════════════════════════════════════════════════════════════

\documentclass[11pt, a4paper]{article}

\usepackage[utf8]{inputenc}
\usepackage[T1]{fontenc}
\usepackage[ngerman]{babel}
\usepackage{helvet}
\renewcommand{\familydefault}{\sfdefault}

\usepackage[margin=1.5cm]{geometry}
\usepackage{array}
\usepackage{booktabs}
\usepackage{tabularx}
\usepackage[table]{xcolor}
\usepackage{enumitem}
\usepackage{fancyhdr}
\usepackage{tikz}
\usepackage{tcolorbox}

% Farben
\definecolor{spanisch}{HTML}{c41e3a}
\definecolor{alternativeB}{HTML}{1565c0}
\definecolor{hellgrau}{HTML}{f5f5f5}
\definecolor{success}{HTML}{2a7a4b}
\definecolor{grau}{HTML}{666666}

\pagestyle{fancy}
\fancyhf{}
\lhead{\small\bfseries LH-10-B}
\chead{\small Spanisch Spätbeginner -- DS 10}
\rhead{\small\textcolor{alternativeB}{\textbf{VARIANTE B: SPEED-DATING}}}
\cfoot{}

\setlength{\parindent}{0pt}
\setlength{\parskip}{0.3em}

\begin{document}

\begin{center}
{\Large\bfseries\color{alternativeB} Speed-Dating: ¿Quién eres?}

\vspace{0.2cm}

{\small Beobachtete Unterrichtsstunde -- 45 Min}
\end{center}

\vspace{0.3cm}

% KURZINFO
\begin{tcolorbox}[colback=alternativeB!10, colframe=alternativeB, boxrule=1pt, arc=0pt]
\textbf{Konzept:} Kurzinterviews (3 Min) mit Partnerwechsel. Timer steuert alles.\\
\textbf{Vorteil:} Timing-sicher, hohe Sprechaktivität, alle sprechen mit mehreren Partnern.\\
\textbf{Materialien:} Fragekarten (8 Stück), Sammelkarten (12 Stück), Timer/Gong
\end{tcolorbox}

\vspace{0.3cm}

% VERLAUFSPLAN
{\large\bfseries Verlaufsplan (45 Min)}

\vspace{0.2cm}

\renewcommand{\arraystretch}{1.3}
\begin{tabularx}{\textwidth}{|c|l|X|c|}
\hline
\rowcolor{alternativeB!20}
\textbf{Min} & \textbf{Phase} & \textbf{Aktivität} & \textbf{SF} \\
\hline
0--5 & Orientierung & Regeln erklären, Sitzordnung (2 Reihen), Fragekarten verteilen & PL \\
\hline
\rowcolor{hellgrau}
5--8 & Runde 1 & Timer 3 Min, A+B stellen sich Fragen & PA \\
\hline
8--9 & Wechsel & \textcolor{spanisch}{\textit{¡Cambio!}} -- B-Reihe rückt 1 Platz weiter & -- \\
\hline
\rowcolor{hellgrau}
9--12 & Runde 2 & Timer 3 Min & PA \\
\hline
12--13 & Wechsel & \textcolor{spanisch}{\textit{¡Cambio!}} & -- \\
\hline
\rowcolor{hellgrau}
13--16 & Runde 3 & Timer 3 Min & PA \\
\hline
16--17 & Wechsel & \textcolor{spanisch}{\textit{¡Cambio!}} & -- \\
\hline
\rowcolor{hellgrau}
17--20 & Runde 4 & Timer 3 Min & PA \\
\hline
20--21 & Wechsel & \textcolor{spanisch}{\textit{¡Cambio!}} & -- \\
\hline
\rowcolor{hellgrau}
21--24 & Runde 5 & Timer 3 Min & PA \\
\hline
24--25 & Wechsel & \textcolor{spanisch}{\textit{¡Cambio!}} & -- \\
\hline
\rowcolor{hellgrau}
25--28 & Runde 6 & Timer 3 Min (letzte Runde) & PA \\
\hline
29--34 & Notieren & Sammelkarte ausfüllen (5 Min): Namen + 2-3 Fakten & EA \\
\hline
34--42 & Reflexion & 4-5 SuS berichten: \textcolor{spanisch}{\textit{¿Qué has aprendido sobre...?}} & PL \\
\hline
42--45 & Cierre & Zusammenfassung, Feedback & PL \\
\hline
\end{tabularx}

\vspace{0.4cm}

% SITZORDNUNG
\begin{minipage}[t]{0.48\textwidth}
{\bfseries Sitzordnung:}

\begin{center}
\begin{tikzpicture}[scale=0.7]
\draw[thick] (0,0) -- (8,0);
\foreach \x in {1,2.5,4,5.5,7} {
    \fill[spanisch] (\x,0.4) circle (0.25);
    \node[white, font=\tiny\bfseries] at (\x,0.4) {A};
}
\foreach \x in {1,2.5,4,5.5,7} {
    \fill[alternativeB] (\x,-0.4) circle (0.25);
    \node[white, font=\tiny\bfseries] at (\x,-0.4) {B};
}
\draw[->, thick] (7.5,-0.4) arc (-90:90:0.5);
\node[right, font=\scriptsize] at (8,0) {B rückt};
\end{tikzpicture}
\end{center}

Zwei Reihen gegenüber. Nach jedem Gong rückt B einen Platz weiter.
\end{minipage}
\hfill
\begin{minipage}[t]{0.48\textwidth}
{\bfseries Differenzierung:}

\textcolor{spanisch}{\textbf{Gruppe A:}} Fragen 1-3 nutzen

\textcolor{alternativeB}{\textbf{Gruppe B:}} Alle 4 Fragen, stellt Rückfragen, passt Tempo an

\vspace{0.3cm}

{\bfseries Fragen auf Karten:}
\begin{enumerate}[nosep, leftmargin=1.2em]
\item \textcolor{spanisch}{\textit{¿Cómo te llamas?}}
\item \textcolor{spanisch}{\textit{¿Cómo eres?}}
\item \textcolor{spanisch}{\textit{¿Qué te gusta hacer?}}
\item \textcolor{spanisch}{\textit{¿Qué profesión quieres?}} {\scriptsize(optional)}
\end{enumerate}
\end{minipage}

\vspace{0.4cm}

% TIPPS
\begin{tcolorbox}[colback=success!10, colframe=success, boxrule=1pt, arc=0pt, title={\textbf{Tipps für die Beobachtung}}]
\begin{itemize}[nosep, leftmargin=1em]
\item \textbf{Timer sichtbar:} Projektor oder großer Timer zeigt verbleibende Zeit
\item \textbf{Gong deutlich:} Klares akustisches Signal (Glocke, Handy-Timer)
\item \textbf{Kein Notieren während Gespräch:} Erst danach auf Sammelkarte!
\item \textbf{Bei Stille:} Selbst Fragen stellen, Impulse geben
\end{itemize}
\end{tcolorbox}

\vspace{0.3cm}

{\small\color{grau}\textbf{Risiko:} NIEDRIG -- Timer steuert alles, keine Abhängigkeit von Schreibgeschwindigkeit.}

\end{document}
